% ============================================================
% Project : THE PHILOSOPHY OF VIABILITY
% Book    : Book I — Ontology of Viability
% Book ID : VS.B01
% File    : section_09_logical_dependency_and_inference_structure.tex
% Section : Chapter 07, Section 09
% Purpose : To examine the logical dependencies among
%           foundational commitments, definitions, axioms,
%           propositions, and derived claims, and to establish
%           standards for representing and evaluating the
%           inferential architecture of the reconstruction.
% ============================================================

\section{Logical Dependency and Inference Structure}
\label{sec:logical-dependency-and-inference-structure}

\subsection{The Architecture of Dependence}

A formal reconstruction is not merely a collection of statements
that happen to be compatible with one another. Its claims occupy
different positions within an organized structure of dependence.
Some claims establish the conceptual conditions under which other
claims can be formulated. Some provide premises for arguments.
Others follow from previously adopted commitments, definitions,
or additional assumptions. Still others remain hypotheses whose
consequences can be explored without treating their premises as
established.

The philosophical significance of this organization is substantial.
A claim may appear convincing when considered in isolation while
depending on an assumption that has not been justified. Conversely,
a proposition may be well supported even though its justification
is distributed across several earlier arguments rather than stated
in the same passage. Without an explicit account of dependency,
the reconstruction risks concealing the structure through which
its claims acquire their support.

The central thesis of this section is that the justification of
an ontological claim must be assessed partly through its position
within a network of dependencies. That position includes the
claim's premises, definitions, inferential rules, interpretive
conditions, and any additional assumptions required for its
establishment.

Logical dependency must therefore be made explicit without
presuming that every form of philosophical dependence is logical
dependence. The reconstruction requires a differentiated account
of the relations among its claims rather than a single undivided
hierarchy.

\subsection{What Counts as a Dependency?}

A dependency exists when the formulation, interpretation,
justification, or derivation of one claim requires another claim
or condition in a specified way.

The term \emph{dependency} is therefore broader than formal
derivability. A proposition may depend on a definition for its
meaning without being logically derived from that definition.
An argument may depend on a methodological convention without
that convention functioning as a premise in the formal proof.
A formal expression may depend on an interpretation that assigns
meaning to its symbols, even though that interpretation is not
itself represented as a formula within the calculus.

At least five kinds of dependency must be distinguished.

\paragraph{Logical dependency.}

A claim \(\varphi\) is logically dependent on a set of premises
\(\Gamma\) when its derivation requires those premises under
specified inference rules:

\[
    \Gamma \vdash_{\mathcal{L}} \varphi.
\]

This relation concerns formal inferential support. It does not
necessarily describe the historical order in which the claims
were formulated or the conceptual order in which they were
introduced.

\paragraph{Definitional dependency.}

A claim may depend on the meaning assigned to one or more of
its terms. If an expression is introduced through a definition,
the interpretation of propositions containing that expression
must respect the definition's stated conditions.

Definitional dependency does not automatically constitute
substantive ontological support. A definition can determine
how a term is used within an argument, but it cannot establish
the reality of what the term purports to designate.

\paragraph{Interpretive dependency.}

A formal expression depends on an interpretation that specifies
the intended domain, the reference of its constants, the
extension of its predicates, and the meaning of its relations
and operators.

A derivation may be valid independently of a particular
interpretation, but a claim about what the derivation means
ontologically requires an adequate interpretive account.

\paragraph{Justificatory dependency.}

A proposition may depend on an argument, evidence, or previously
defended commitment for its acceptance. This relation concerns
the grounds on which a claim is warranted, rather than merely
the formal possibility of deriving it.

Justificatory dependency may include logical derivation, but it
may also involve explanatory arguments, comparative analysis,
or the evaluation of alternative interpretations.

\paragraph{Methodological dependency.}

A claim may depend on a methodological decision concerning the
scope, level, or procedure of analysis. For example, a formal
comparison may require that two expressions be interpreted
within a common domain or that a particular respect of
transformation be held fixed.

Such a decision can be necessary for the inquiry without
constituting an ontological premise.

These forms of dependency may overlap. A proposition can be
logically derived from premises whose interpretation depends
on a definition, while its ontological acceptance also depends
on an independent justificatory argument. The relations must
therefore be identified according to their actual roles rather
than collapsed into one generic relation of dependence.

\subsection{The Distinction Between Presentation Order and
Dependency Order}

The order in which claims appear in a book is not necessarily
the order in which they depend on one another.

A chapter may introduce a concept before presenting the
argument that justifies its use. A proposition may be stated
before its formal derivation for explanatory reasons. A later
section may clarify a definition that was provisionally used
in an earlier discussion. None of these arrangements necessarily
indicates a logical inconsistency.

The distinction between presentation order and dependency
order is essential.

\begin{itemize}
    \item \textbf{Presentation order} concerns the sequence in
    which concepts and arguments are introduced to the reader.

    \item \textbf{Conceptual order} concerns the relations of
    meaning and intelligibility among concepts.

    \item \textbf{Justificatory order} concerns the sequence in
    which reasons and grounds must be supplied for accepting
    particular claims.

    \item \textbf{Logical order} concerns the relations through
    which conclusions follow from premises under specified
    rules.
\end{itemize}

These orders may coincide in some cases, but no general identity
among them should be assumed.

For example, the concept of identity may be introduced before
a complete formal account of its relation to being is available.
This presentation decision does not establish that identity
is logically prior to being, or that being is conceptually
dependent on identity. It merely reflects an explanatory choice
made in arranging the discussion.

Similarly, a proposition introduced early in the exposition
may receive its full justification only after several
intermediate distinctions have been established. Its position
in the text does not determine the direction of its logical
dependence.

\subsection{A Typed Structure of Claims}

The dependency structure must preserve the distinctions among
the different functional forms of claims.

The principal categories include:

\begin{enumerate}
    \item \textbf{Foundational commitments}, which express
    adopted ontological orientations or constraints.

    \item \textbf{Definitions}, which establish the meanings
    and conditions of use of specified expressions.

    \item \textbf{Axioms}, which serve as starting expressions
    within a formal calculus.

    \item \textbf{Propositions}, which state claims that may
    require argument or derivation.

    \item \textbf{Theses}, which articulate substantive positions
    defended through a broader philosophical argument.

    \item \textbf{Hypotheses}, which introduce provisional
    assumptions for examination.

    \item \textbf{Arguments}, which provide inferential or
    explanatory support for conclusions.

    \item \textbf{Open problems}, which identify questions
    whose resolution has not been established.
\end{enumerate}

These categories describe the functional or logical form of
claims and argumentative elements. They must not be confused
with their epistemic or justificatory status.

A proposition may be conditionally supported. A thesis may
remain provisional. A definition may be stipulated for a
restricted formal purpose. An axiom may be accepted within a
calculus while its ontological adequacy remains open to
examination.

The dependency structure should therefore record two independent
dimensions:

\[
    \text{Claim Classification}
    =
    \text{Functional Form}
    +
    \text{Epistemic Status}.
\]

This expression is a classification principle rather than a
mathematical identity. It indicates that the type of a claim
and the grounds for accepting it answer different questions.

A dependency map that records only the formal type of each
node would be incomplete. It must also preserve the status
and conditions under which the node is being used.

\subsection{Foundational Commitments and Their Dependencies}

The three foundational principles established in the preceding
sections provide a useful example of how dependencies must
be represented without imposing an unjustified hierarchy.

Let:

\[
    AD = \text{Analytical Distinction},
\]

\[
    OI = \text{Ontological Inseparability},
\]

\[
    AP = \text{Absence of Ontological Priority}.
\]

Their joint adoption may be represented as:

\[
    F_{BI} := AD \land OI \land AP.
\]

This representation records the joint presence of the three
commitments within the adopted framework. It does not establish
a derivation of one principle from another.

The dependency structure must therefore preserve the following
features:

\begin{enumerate}
    \item \(AD\), \(OI\), and \(AP\) are distinct commitments.

    \item Their joint adoption does not make them logically
    interchangeable.

    \item The conjunction \(F_{BI}\) does not establish that
    the three principles are exhaustive of every possible
    ontological commitment concerning being and identity.

    \item Any proposition derived from their conjunction must
    identify the specific premises and interpretive conditions
    that support it.

    \item No additional ontological principle may be introduced
    merely because it would make the formal system more
    symmetrical or convenient.
\end{enumerate}

The principles thus function as jointly adopted foundations
without being reduced to a single undifferentiated axiom.
Their relation is one of coordinated commitment, not an
established relation of mutual derivability.

\subsection{Dependency Analysis of Analytical Distinction}

The formal expression

\[
    q_B \neq q_I
\]

records the distinction between the analytical questions or
roles associated with being and identity.

A proposition that depends on this expression may legitimately
claim that the analysis must preserve the difference between
the two questions. It may also use the distinction to organize
definitions, arguments, and formal representations.

However, a further claim that being and identity are two
independently existing entities does not follow from this
expression alone.

Such a conclusion would require an additional premise connecting
analytical differentiation with ontological separation. The
presence of that premise would need to be justified independently,
and it would have to be examined in relation to the commitment
to ontological inseparability.

The dependency structure should therefore distinguish:

\[
    \text{Analytical Differentiation}
\]

from

\[
    \text{Ontological Separation}.
\]

The first is supported by the adopted analytical distinction.
The second is not licensed by that distinction and cannot be
introduced as its automatic consequence.

This example illustrates why dependency analysis must preserve
the semantic content of premises. A formal relation between
symbols cannot be extended to a different ontological relation
without an explicit and defensible bridge.

\subsection{Dependency Analysis of Ontological Inseparability}

The co-applicability schema introduced in the preceding section
is:

\[
    \forall x \in D_E
    \bigl(E(x) \rightarrow (B(x) \land I(x))\bigr).
\]

Within the specified domain and interpretation, the schema
supports the conclusion that an existent satisfying \(E(x)\)
is represented as satisfying both \(B(x)\) and \(I(x)\).

The inferential dependency can be stated as follows:

\[
\begin{aligned}
    &\forall x \in D_E
    \bigl(E(x) \rightarrow (B(x) \land I(x))\bigr),\\
    &E(a),\\
    &a \in D_E\\
    &\therefore B(a) \land I(a).
\end{aligned}
\]

The conclusion depends on the universal schema, the individual
instance, and the admissibility of \(a\) within the specified
domain.

It does not follow that the schema captures every aspect of
ontological inseparability. Nor does it establish that the
predicates \(B\) and \(I\) have identical meanings or extensions
under every possible interpretation.

The dependency structure must therefore distinguish the
co-applicability result from the broader ontological commitment
it is intended to represent.

A formal system may use the schema as a partial representation
of the commitment, but the stronger interpretation remains
dependent on an ontological account that explains why being
and identity are inseparable without being analytically
indistinguishable.

This distinction also prevents an invalid dependency from
being introduced between ontological inseparability and
predicate equivalence.

\subsection{Dependency Analysis of the Absence of Ontological
Priority}

The formal representation of the absence of ontological priority
is:

\[
    \neg(B \prec_O I)
    \land
    \neg(I \prec_O B).
\]

The expression depends on the definition of the relation
\(\prec_O\). Unless the meaning and conditions of that relation
are specified, the formal negations remain insufficiently
interpreted.

Once the relation is adequately specified, conjunction
elimination permits the two component claims to be considered
separately:

\[
    \neg(B \prec_O I)
\]

and

\[
    \neg(I \prec_O B).
\]

The result concerns the two directions of the specified
ontological priority relation. It does not establish the
absence of every other relation between being and identity.

In particular, the following relations require independent
analysis:

\begin{itemize}
    \item conceptual dependence;
    \item explanatory ordering;
    \item temporal succession;
    \item causal relation;
    \item epistemic dependence;
    \item methodological ordering.
\end{itemize}

No one of these relations should be treated as equivalent to
ontological priority without an explicit argument.

Consequently, a proposition concerning the absence of
conceptual dependence, for example, cannot be derived merely
from the denial of ontological priority. Its justification
would require premises addressing conceptual dependence
directly.

\subsection{Dependency Chains and Intermediate Results}

Many philosophical arguments cannot be represented adequately
as a direct transition from foundational commitments to a
final conclusion. They require intermediate propositions,
definitions, and interpretive clarifications.

A dependency chain may be represented schematically as:

\[
    \Gamma_0
    \longrightarrow
    \Gamma_1
    \longrightarrow
    \Gamma_2
    \longrightarrow
    \varphi,
\]

where each transition indicates a justified inferential or
argumentative relation rather than a mere sequence of presentation.

For such a chain to be reliable, each intermediate stage must
preserve the conditions inherited from the preceding stages.

Suppose that:

\[
    \Gamma_0 \vdash_{\mathcal{L}} \psi_1,
\]

and

\[
    \Gamma_1, \psi_1
    \vdash_{\mathcal{L}} \psi_2.
\]

A subsequent derivation may establish:

\[
    \Gamma_2, \psi_2
    \vdash_{\mathcal{L}} \varphi.
\]

The conclusion \(\varphi\) depends on the complete set of
assumptions required by the chain. It cannot be presented
as following from \(\Gamma_0\) alone unless the additional
premises in \(\Gamma_1\) and \(\Gamma_2\) are themselves
derived from \(\Gamma_0\) or otherwise independently justified.

This point is essential when a conclusion appears several
sections after the concepts on which it depends. The distance
between a premise and a conclusion in the text does not
eliminate their inferential relationship, but neither does
their proximity establish one.

\subsection{Direct, Indirect, and Conditional Dependencies}

For analytical purposes, dependencies should be distinguished
according to their structure.

\paragraph{Direct dependency.}

A claim directly depends on another when the latter is used
as a premise, definition, interpretive condition, or explicit
justificatory ground in establishing the former.

\paragraph{Indirect dependency.}

A claim indirectly depends on another when the dependency
is mediated through one or more intermediate claims.

For example:

\[
    \alpha \rightarrow \beta,
    \qquad
    \beta \rightarrow \gamma.
\]

If both transitions are adequately established, \(\gamma\)
may depend indirectly on \(\alpha\).

\paragraph{Conditional dependency.}

A claim is conditionally dependent when its support requires
an assumption that has not been adopted unconditionally.

If:

\[
    \Gamma \cup \{\delta\}
    \vdash_{\mathcal{L}} \varphi,
\]

then \(\varphi\) remains dependent on \(\delta\), unless
\(\delta\) is independently established or eliminated through
a separate derivation.

\paragraph{Interpretive dependency.}

A claim may depend on a specific interpretation of its
formal expressions. A change in that interpretation may
alter the meaning or applicability of the conclusion even
if the symbolic derivation remains unchanged.

\paragraph{Definitional dependency.}

A proposition may depend on the conditions stipulated by
one or more definitions. If those definitions are revised,
the proposition may require reinterpretation or reconstruction.

These categories are not mutually exclusive. A proposition
may be directly dependent on a conditional premise whose
interpretation is itself dependent on a definition.

\subsection{Dependency and Non-Circular Justification}

A dependency structure becomes circular when a claim is
presented as justified by another claim whose justification
already presupposes the first claim, without an independent
ground that breaks the circle.

A simple circular pattern can be represented as:

\[
    \alpha \text{ is justified by } \beta,
\]

\[
    \beta \text{ is justified by } \alpha.
\]

Such a pattern does not necessarily produce a formal contradiction.
Its defect is justificatory: neither claim receives independent
support from the other.

The presence of reciprocal relations in an ontological theory
does not automatically constitute circular reasoning. Two
features may be mutually related or jointly necessary without
one being used as the sole justification of the other.

The three foundational principles provide an important example.
Analytical distinction, ontological inseparability, and absence
of ontological priority are jointly adopted commitments. Their
coordination does not require that one be proved exclusively
from another.

A circularity arises only if the argument claims to establish
one principle by relying on the very content of that principle
as an unstated premise.

To evaluate possible circularity, the reconstruction must ask:

\begin{enumerate}
    \item Is the conclusion already assumed, explicitly or
    implicitly, among the premises?

    \item Does an intermediate definition presuppose the
    claim that the argument is intended to establish?

    \item Does the argument move between different meanings
    of a term in a way that makes the conclusion appear
    necessary?

    \item Is a reciprocal ontological relation being confused
    with a reciprocal justificatory relation?

    \item Can the argument be reconstructed with its premises
    and inferential steps made explicit?
\end{enumerate}

The objective is not to eliminate every form of conceptual
reciprocity. It is to prevent reciprocal relations from
concealing the absence of independent justification.

\subsection{Dependency, Redundancy, and Independence}

The dependency structure also provides a basis for examining
whether claims are redundant or independent.

Two claims may be related in several ways:

\begin{itemize}
    \item one may be derivable from the other;
    \item both may be derivable from a common set of premises;
    \item one may clarify the meaning of the other without
    being logically equivalent to it;
    \item they may overlap in content while retaining distinct
    justificatory roles;
    \item neither may be derivable from the other within the
    adopted system.
\end{itemize}

Redundancy should not be inferred merely from thematic
similarity. Two principles may concern the same subject while
performing different analytical functions.

Likewise, independence cannot be established merely because
no derivation has yet been presented. A claim may be
independent relative to a specified formal system, but
demonstrating such independence requires an appropriate
method, such as constructing interpretations or models in
which one claim holds while the other does not.

This distinction is important for the three foundational
principles. Their conceptual and justificatory separation
must be preserved. Whether particular formal formulations
are independent within a chosen calculus is a further
technical question and cannot be settled solely by their
different names or verbal descriptions.

\subsection{Dependency Under Revision}

A philosophical reconstruction remains open to correction.
Definitions may be refined, hypotheses may be rejected,
and formal representations may be replaced. Such revisions
can affect claims that depend on the revised element.

The consequences of revision depend on the kind of dependency
involved.

If a definition changes, propositions using the defined
expression may require reinterpretation.

If a premise is withdrawn, every conclusion whose derivation
requires that premise must be reconsidered.

If an inference rule is modified, derivations relying on that
rule must be re-evaluated.

If an interpretation is replaced, the ontological significance
of formal expressions may change even when their symbolic
relations remain intact.

If a hypothesis is rejected, conditional consequences derived
from it lose their support as consequences of an accepted
premise. They may remain useful as descriptions of what would
follow under the rejected hypothesis, but they cannot retain
the status of established conclusions.

Revision must therefore be propagated through the dependency
structure. Correcting one claim while leaving its dependent
claims untouched can create hidden inconsistencies in the
reconstruction.

This requirement does not mean that every revision invalidates
the entire system. A well-organized dependency structure makes
it possible to identify precisely which conclusions are
affected and which remain independently supported.

\subsection{A Practical Protocol for Dependency Analysis}

Every substantial proposition introduced in the formal
reconstruction should be examined according to a consistent
protocol.

\begin{enumerate}
    \item \textbf{Identify the claim.}
    State the proposition or commitment in a form sufficiently
    precise for its dependencies to be examined.

    \item \textbf{Classify its function.}
    Identify whether it is a foundational commitment, definition,
    axiom, proposition, thesis, hypothesis, argument, or open
    problem.

    \item \textbf{Record its epistemic status.}
    Specify whether it is foundationally adopted, stipulated,
    argued or derived, conditionally supported, provisional,
    hypothetical, unresolved, or rejected.

    \item \textbf{Identify direct premises.}
    Record every premise required for its derivation or
    justification.

    \item \textbf{Identify definitions and interpretations.}
    Specify the meanings and interpretive conditions on which
    the claim depends.

    \item \textbf{Reconstruct intermediate steps.}
    Make indirect dependencies and essential transitions
    explicit.

    \item \textbf{Check inferential validity.}
    Examine whether the conclusion follows under the adopted
    rules and assumptions.

    \item \textbf{Check scope and status preservation.}
    Ensure that the conclusion does not exceed the domain,
    conditions, or justificatory status of its supporting
    premises.

    \item \textbf{Test for circularity and redundancy.}
    Determine whether the argument presupposes its conclusion
    or repeats an already established claim without adding
    explanatory or inferential value.

    \item \textbf{Record unresolved dependencies.}
    Identify any premise, definition, or interpretive condition
    that remains provisional or requires further examination.
\end{enumerate}

This protocol is intended to support both the construction
and the later audit of the formal system. It does not require
every philosophical claim to be reduced to symbolic notation.
It requires that the grounds and dependencies of substantial
claims remain traceable.

\subsection{The Limits of a Dependency Graph}

A dependency graph can be useful for representing relations
among claims, but it is not a substitute for their philosophical
analysis.

A graph may represent claims as nodes and selected dependency
relations as directed edges. Such a representation can reveal
chains of inference, shared premises, isolated claims,
conditional branches, and possible points of failure.

However, a graph can also mislead if its structure is not
carefully defined.

First, an arrow must have a specified meaning. It may indicate
logical derivation, definitional dependence, interpretive
dependence, or justificatory support. These relations must
not be represented by an undifferentiated arrow if doing so
would conceal their differences.

Second, the direction of an edge must reflect the actual
dependency relation. The order in which two claims appear
in the book is not sufficient to determine the direction.

Third, a graph may simplify complex dependencies involving
several premises, background assumptions, and interpretive
conditions. A single edge can conceal a conjunction of
requirements.

Fourth, a graph cannot by itself establish the validity of
an inference. It can display a proposed dependency, but the
legitimacy of that dependency must be examined through the
relevant argument or derivation.

Fifth, a graph may create the appearance of a fixed hierarchy
where the underlying relations are conditional, reciprocal
in a non-justificatory sense, or dependent on alternative
interpretations.

For these reasons, any future visual representation of the
dependency structure must distinguish the types of relations
being displayed and must be accompanied by an explanation
of its conventions. A graphical representation is an aid
to analysis, not an independent source of ontological or
logical warrant.

\subsection{Open Problems in the Architecture of Dependence}

The present account establishes a framework for examining
logical and justificatory dependencies, but several questions
remain open.

First, the complete set of foundational commitments required
by the ontology has not been conclusively determined. The
three principles concerning being and identity provide a
central orientation, but their sufficiency for every later
ontological question must not be presumed.

Second, the precise formal relations between definitions,
axioms, propositions, and interpretive conditions require
further specification. Not every conceptual dependency can
be represented within a single formal calculus.

Third, the independence of particular principles and
propositions requires technical examination within an
explicitly specified system. Conceptual distinction and
formal independence are related but non-identical questions.

Fourth, the conditions under which different formal
interpretations preserve the same ontological commitments
remain to be determined.

Fifth, a complete dependency audit must examine whether
revisions to definitions or foundational assumptions affect
previously formulated claims elsewhere in the book.

Sixth, the relation between explanatory support and formal
derivability requires continued attention. Some ontological
arguments may be philosophically important without being
deductive proofs, and their justificatory role must be
represented accordingly.

These questions define the limits of the present section.
They do not prevent the reconstruction from proceeding.
Rather, they establish the conditions under which its
subsequent formal claims can be organized, assessed, and
revised without concealing their dependencies.

\subsection{Concluding Principle}

A formal ontological reconstruction acquires coherence not
merely by assembling compatible statements, but by making
their relations of dependence explicit and assessable.

The governing principle is:

\begin{quote}
Every substantial claim must remain traceable to the
definitions, commitments, premises, interpretations, and
inferential procedures on which its formulation or
justification depends. No claim may acquire a stronger
status or a wider scope than its dependencies warrant.
\end{quote}

This principle establishes a disciplined architecture for
the reconstruction. It permits foundational commitments
to guide formal development while preserving the distinction
between adopted premises, stipulated meanings, derived
consequences, conditional results, and unresolved questions.

The next section examines the consistency, independence,
and possible conflicts among the claims organized within
this architecture.

% ============================================================
% END OF SECTION 09 OF CHAPTER 07
% ============================================================